\chapter{\label{chap:conclusion}Conclusion and Future Work}

\section{Summary of Contributions}

This thesis develops Twin2MultiCloud as a bounded research proof of concept for
operationalizing a layered, cost-aware Digital Twin model across AWS, Azure,
and Google Cloud. The central contribution is the traceable connection between
typed workload intent, functionally admissible provider bundles, evidence-based
cost calculation, an exact deployment specification, supervised
infrastructure provisioning, runtime verification, and cleanup.

The implementation contributes a standalone Six-layer Eventing contract, a
functional-completeness gate, frozen pricing and formula evidence, a cost-only
selection strategy, digest-linked RTA/RDS/Manifest contracts, graph-derived
credential readiness, immutable Twin lifecycle, durable deployment operations,
and a unified Management API and Flutter workflow. The predecessor optimizer
and Deployer are retained as architectural lineage, while obsolete product
features are not part of the executable research boundary.

% Complete the following paragraph only after the supervised evidence is
% frozen: summarize the RQ1 operationalization result, RQ2 comparability
% boundaries, RQ3/RQ3.1 cost differences, RQ3.2 Eventing effect, provider
% blockers, observed-vs-estimated cost, and residual cleanup result. Cite the
% exact scenario and digest records; do not insert offline fixtures as live
% outcomes.

\section{Future Directions}
\label{sec:future-work}

The following extensions are research topics rather than a product roadmap.

\subsection{Additional Optimization Objectives}

The internal scoring seam could support future latency, sustainability,
resilience, or policy studies. Each objective would require comparable metrics,
measurement provenance, normalization, trade-off semantics, and an independent
evaluation. Merely adding disabled objectives or weighted configuration fields
would not demonstrate the capability.

\subsection{Alternative Architecture Contracts}

The PoC executes only \texttt{six-layer-eventing@1}. Alternative Digital Twin
decompositions could be compared through separately versioned contracts,
admissibility rules, and evaluation matrices. They should not inherit
implicitly from the present profile or appear as unvalidated runtime choices.

\subsection{Least-Privilege Credential Lifecycle}

A production-oriented study could derive reduced deployment principals from
the resolved graph, rotate and revoke them, integrate organization policy, and
formally verify permissions. This would require a dedicated threat model,
provider-specific recovery procedures, and clear operational ownership beyond
the administrator CloudConnections used for isolated thesis accounts.

\subsection{Mutable Deployment Semantics}

Supporting in-place changes would require optimizer invalidation, Terraform
replacement analysis, state migration, rollback, partial-failure recovery,
identity rotation, and a new evidence model. These properties should be
evaluated together rather than represented as an unrestricted Edit action.

\subsection{Broader Empirical Evidence}

Future evaluations could study concurrent deployments, Medium and Large
workloads, long-running delivery and replay behavior, reconciliation against
provider bills, regional variation, additional provider bundles, and
organization-scale account boundaries. These questions are not answered by
the nine supervised Small scenarios of this thesis.
